\section{Combination Sum III}
\label{sec:rec-combination-sum-3}

\problemstatement{Find all valid combinations of exactly $k$
\textbf{distinct} numbers, chosen from $1$ through $9$ (each usable at
most once across all combinations, i.e.\ no repeats within one
combination), that sum to exactly $n$.}

\begin{patternbox}
A further specialization of Section~\ref{sec:rec-combination-sum-2}'s
no-reuse, advancing recursion: the candidate pool is now the fixed,
already-sorted range $1$ through $9$ (so no separate sort step or
duplicate-skip is needed, since this fixed pool has no duplicates by
construction), and one new constraint is added -- the combination's
\textbf{length} must be exactly $k$, not just its sum.
\end{patternbox}

\subsection*{Understanding the problem carefully}

The recursion is the same advancing, no-reuse shape as Combination Sum
II, but simplified by the fixed candidate pool $\{1,\dots,9\}$ (already
sorted, no duplicates to worry about), and with the base case now
checking \emph{two} conditions instead of one: both the running count of
chosen numbers must equal $k$, and the running remaining target must
have reached exactly $0$, at the same time.

\begin{ideabox}[Same Template, with a Length Constraint Added to the Base Case]
$\textsc{Solve}(\textit{start}, \textit{remaining}, \textit{current})$:
if $|\textit{current}| = k$: record $\textit{current}$ if
$\textit{remaining}=0$ (otherwise this length-$k$ combination simply
doesn't sum correctly, and we discard it); return either way (no further
recursion needed once length $k$ is reached). Otherwise, for $i$ from
$\textit{start}$ to $9$: if $i > \textit{remaining}$: break (the same
sorted-range pruning as Combination Sum II). Append $i$; recurse
$\textsc{Solve}(i+1, \textit{remaining}-i, \textit{current})$; remove
(undo).
\end{ideabox}

\subsection*{Why checking both conditions together at length $k$ is necessary}

A combination might reach exactly $k$ chosen numbers without summing to
$n$ (too small or too large), or might sum to exactly $n$ using fewer or
more than $k$ numbers -- neither partial success is acceptable, since the
problem requires \emph{both} conditions simultaneously. Stopping the
recursion the instant the count reaches $k$ (regardless of whether the
sum is currently correct) avoids wasted exploration beyond the required
length, while still checking the sum condition explicitly before
recording, since reaching length $k$ alone does not guarantee the target
sum has also been met.

\subsection*{Algorithm}

\begin{enumerate}
  \item $\textsc{Solve}(\textit{start}, \textit{remaining}, \textit{current})$:
  \begin{enumerate}
    \item If $|\textit{current}|=k$: if $\textit{remaining}=0$: record
          $\textit{current}$; return.
    \item For $i$ from $\textit{start}$ to $9$:
    \begin{itemize}
      \item If $i > \textit{remaining}$: break.
      \item Append $i$; recurse $\textsc{Solve}(i+1, \textit{remaining}-i,
            \textit{current})$; remove (undo).
    \end{itemize}
  \end{enumerate}
  \item Call $\textsc{Solve}(1, n, [\,])$.
\end{enumerate}

\subsection*{Dry run}

\dryrun{Take $k=3$, $n=7$.}

\begin{itemize}
  \item $\textsc{Solve}(1,7,[])$: $i=1$: append, recurse
        $\textsc{Solve}(2,6,[1])$.
  \begin{itemize}
    \item $i=2$: append, recurse $\textsc{Solve}(3,4,[1,2])$.
    \begin{itemize}
      \item $|\textit{current}|=2 \ne 3$ yet. $i=3$: $3\le4$, append,
            recurse $\textsc{Solve}(4,1,[1,2,3])$: $|\textit{current}|=3=k$,
            but $\textit{remaining}=1\ne0$: discard. $i=4$: $4>1$,
            break.
    \end{itemize}
    \item Undo. $i=3$ (at $[1]$ level): append, recurse
          $\textsc{Solve}(4,3,[1,3])$: $i=4$: $4>3$, break (no length-3
          combination found here).
  \end{itemize}
  \item Continue similarly... eventually reaching $i=1,i=2,i=4$:
        $\textsc{Solve}(2,6,[1])\to i=4$: $4\le6$, append, recurse
        $\textsc{Solve}(5,2,[1,4])$: $i=5$: $5>2$, break.
  \item Trying $i=1, i=6$ at the $[1]$ level... continuing the full
        search eventually finds $[1,2,4]$ ($1+2+4=7$, length $3$) as a
        valid recorded combination.
\end{itemize}

(The complete search also finds no other length-$3$ combinations from
$\{1,\dots,9\}$ summing to $7$ besides $[1,2,4]$, which the C++ code
below confirms by direct execution.)

\subsection*{C++ implementation}

\begin{lstlisting}
#include <bits/stdc++.h>
using namespace std;

void solve(int start, int remaining, int k, vector<int>& current,
           vector<vector<int>>& result) {
    if ((int)current.size() == k) {
        if (remaining == 0) result.push_back(current);
        return;
    }
    for (int i = start; i <= 9; i++) {
        if (i > remaining) break;

        current.push_back(i);
        solve(i + 1, remaining - i, k, current, result);
        current.pop_back();
    }
}

vector<vector<int>> combinationSum3(int k, int n) {
    vector<vector<int>> result;
    vector<int> current;
    solve(1, n, k, current, result);
    return result;
}

int main() {
    for (auto& comb : combinationSum3(3, 7)) {
        cout << "[ ";
        for (int x : comb) cout << x << " ";
        cout << "] ";
    }
    cout << endl;
    return 0;
}
\end{lstlisting}

\begin{complexitybox}
\textbf{Time Complexity.} The candidate pool is fixed at $9$ values, so
the search space is bounded by $\binom{9}{k}$, a constant independent of
any input size -- effectively $O(1)$ for this fixed-range problem, though
conventionally still described as exponential in the (here, fixed)
pool size.
\textbf{Space Complexity.} $O(k)$ for the recursion depth.
\end{complexitybox}

\begin{takeawaybox}
A fixed, small, already-sorted, duplicate-free candidate pool (like
$1$--$9$ here) simplifies a backtracking template by removing the need
for an initial sort or a same-level duplicate-skip check -- but the core
advancing, no-reuse recursion, plus sorted-pool pruning on the
remaining target, is unchanged from Combination Sum II. Adding a second
base-case condition (here, an exact length $k$) is a small, composable
extension worth recognizing as its own reusable idea.
\end{takeawaybox}
